\documentclass{situationist-preprint}

\title{Truths About Absences Without Absences}
\msid{MET-16-a}
\msfamily{MET-16}
\msdiscipline{Metaphysics}
\mstype{Principal manuscript}
\mskeywords{absences; negative truths; truthmakers; truth supervenes on being; causation by absence; holes}

\begin{document}
\maketitle

\begin{abstract}
We say that the absence of rain caused the drought, that Pierre's absence from the café was palpable, and that there is no beer in the fridge. Are there absences, entities that these sentences are about? We argue that there are not, but that there are true sentences about absences, and that this suffices for every purpose for which absences have been posited. The argument has two parts. First, negative truths do not require truthmakers: they are true because of what there is and what there is not, in the sense captured by the principle that truth supervenes on being. Secondly, the remaining roles of absences, as causes and as objects of perception, are discharged by counterfactual dependence and by the presence of perceptible things, without any entity that is an absence. The view is distinguished from the view that holes exist, which is compatible with it: holes are dependent particulars that can be filled and have hosts, and are not absences.
\end{abstract}

\section{The problem}

Absences figure in true sentences of several kinds. There are negative existentials (\enquote{there is no beer in the fridge}), statements of causation by absence (\enquote{the gardener's failure to water the plants caused them to die}), and reports of perception (\enquote{Pierre's absence was what I noticed when I came into the café}). The last is Sartre's example: entering a café to meet Pierre, who is not there, one sees his absence, which organizes the whole scene \citep{Sartre1943}. If these sentences are true, and if truth requires that something in the world make the sentence true, it seems that there must be absences.

The literature offers three responses. \citet{Kukso2006} argues that absences are real, since they are perceived and cause things. \citet{Armstrong2004} denies that there are absences as such, but holds that negative truths require truthmakers, and supplies them in the form of totality facts: the fact that these are all the things in the fridge makes it true that there is no beer there. \citet{Lewis2001} and \citet{Bigelow1988} deny that negative truths require truthmakers at all: a negative existential is true because nothing makes its negation true. We defend the third response, and extend it to causation and perception.

\section{Framework}

\begin{definition}\label{d:absence}
An \emph{absence} would be an entity whose nature is to be the non-existence or non-occurrence of something of a given kind at a place or time.
\end{definition}

\begin{principle}[Truth supervenes on being]\label{p:tsb}
There is no difference in which propositions are true without a difference in what things there are or in which properties and relations they instantiate.
\end{principle}

\begin{principle}[No idle entities]\label{p:idle}
We should not posit entities of a kind if every truth whose explanation is the reason for positing them is explained without them.
\end{principle}

Principle~\ref{p:tsb} is the formulation of the truthmaker intuition defended by \citet{Bigelow1988} and \citet{Lewis2001}. It is weaker than truthmaker maximalism, the thesis that every truth has a truthmaker. It requires that two worlds differing in some truth differ in being; it does not require that the difference be the presence of an extra entity in the world in which the truth holds. A world in which there is no beer in the fridge differs from one in which there is, in that the second contains some beer in the fridge that the first lacks. The difference is a difference in being, but it is not the presence of an absence in the first world.

\section{The argument}

\begin{lemma}\label{l:neg}
Every negative existential truth is explained without absences.
\end{lemma}

\begin{proof}
Let \enquote{there is no $F$ at $r$} be true. By Principle~\ref{p:tsb}, it suffices that every world in which it is false differs in being from ours, and every such world contains an $F$ at $r$, which ours does not. The truth is explained by what there is not, that is, by the non-existence of any $F$ at $r$, and not by the existence of a further thing. The totality facts of \citet{Armstrong2004} are not needed either: they are posited to satisfy truthmaker maximalism, which Principle~\ref{p:tsb} does not require.
\end{proof}

\begin{lemma}\label{l:cause}
Every truth of the form \enquote{the absence of $C$ caused $E$} is explained without absences.
\end{lemma}

\begin{proof}
On the account of causation defended in MET-04, a cause is a counterfactual difference-maker selected by norms of normality. The sentence \enquote{the absence of water caused the plant's death} is true when the plant died, it would not have died had it been watered, and watering is the norm against which the failure stands out. Each of these conditions is stated in terms of the plant, the water and the counterfactual relations between them. None requires an entity that is the absence of water: the counterfactual \enquote{had there been water, the plant would not have died} is about the presence of water in other possible situations, not about an absence in this one \citep{Lewis2004}.
\end{proof}

\begin{lemma}\label{l:perc}
Every truth of the form \enquote{$S$ perceived the absence of $x$} is explained without absences.
\end{lemma}

\begin{proof}
When I see that Pierre is not in the café, I see the café, the people in it and the tables, and I see that none of the people is Pierre. What I perceive is the presence of what is there, under an expectation that Pierre would be among it. The perceptual state is directed at existing things and represents them as not including Pierre. The report \enquote{I saw Pierre's absence} is true because of this state, which requires no entity that is Pierre's absence \citep{Martin1996}.
\end{proof}

\begin{theorem}\label{t:main}
There are no absences, and true sentences about absences require none.
\end{theorem}

\begin{proof}
The truths that would motivate positing absences are negative existentials, causal statements and perceptual reports. By Lemmas~\ref{l:neg}, \ref{l:cause} and~\ref{l:perc}, each of these is explained without absences. By Principle~\ref{p:idle}, we should not posit absences. The sentences about absences remain true, and their truth suffices for every purpose for which absences were posited.
\end{proof}

\section{Absences and holes}

A hole in a piece of cheese is located where there is no cheese, and it is natural to describe it as an absence of cheese. MET-15 argues that holes exist. If holes were absences, the two theses would be inconsistent. They are not, because holes are not absences.

A hole is a dependent immaterial particular located at the boundary of its host. It has a host, whose surface configuration fixes its identity; an absence of cheese from a region has no host. A hole can be filled with water and remain the same hole; an absence of matter from a region ceases when matter arrives there. And the truth that the cheese has a hole is a positive truth about the configuration of the cheese and about the particular that the configuration bounds, while the truth that there is no cheese in a region is a negative existential, true because there is no cheese there. Theorem~\ref{t:main} denies the existence of entities of the second kind and is silent about the first. The coincidence, in an unfilled hole, of the hole's location with a region from which matter is absent does not identify the two.

\section{Objections}

\begin{objection}[Causal efficacy]
Absences make a difference: the drought killed the crops. Whatever makes a difference exists.
\end{objection}

\begin{reply}
Making a difference, on the account of MET-04, is a matter of counterfactual dependence. The crops' death depends counterfactually on whether it rained, and the dependence is a relation between the occurrence of rain in other situations and the death in this one. The principle that whatever makes a difference exists is true when \enquote{whatever makes a difference} ranges over causal relata in the sense of entities that bring about effects by their activity; it is false when it ranges over everything that can be cited in a true causal explanation, which includes the non-occurrence of events \citep{Beebee2004}.
\end{reply}

\begin{objection}[Truthmaker maximalism]
Every truth has a truthmaker. Negative truths therefore require either absences or totality facts.
\end{objection}

\begin{reply}
Truthmaker maximalism is the premise we reject. Its motivation is the thought that truth depends on the world, and Principle~\ref{p:tsb} captures that thought without requiring that each truth be made true by the existence of something. The cost of maximalism is high: it requires either absences, which are idle by Theorem~\ref{t:main}, or totality facts, which are negative facts in disguise \citep{Molnar2000}.
\end{reply}

\begin{objection}[Reference]
\enquote{Pierre's absence} is a singular term, and it occurs in true sentences. A singular term in a true atomic sentence refers.
\end{objection}

\begin{reply}
The sentences in which it occurs are not atomic in their logical form. \enquote{Pierre's absence was noticed} is true because someone noticed that Pierre was not there, and the nominalization is a device for speaking compactly about a negative fact, in the way that \enquote{the average family} is a device for speaking about families. The quantification over absences in ordinary speech, unlike the quantification over holes, admits a paraphrase that preserves every inference we draw from it.
\end{reply}

\section{Conclusion}

There are no absences. Negative existentials are true because of what there is not, causal statements about absences are true because of counterfactual dependence, and perceptual reports of absences are true because of what is perceived to be present. The sentences about absences are true, and their truth suffices. Holes, which do exist, are not absences, and the two theses are consistent.

\begin{thebibliography}{9}\small
\bibitem[Armstrong(2004)]{Armstrong2004} Armstrong, D.~M. (2004). \emph{Truth and Truthmakers}. Cambridge University Press.
\bibitem[Beebee(2004)]{Beebee2004} Beebee, H. (2004). Causing and nothingness. In J. Collins, N. Hall and L.~A. Paul (Eds.), \emph{Causation and Counterfactuals}. MIT Press.
\bibitem[Bigelow(1988)]{Bigelow1988} Bigelow, J. (1988). \emph{The Reality of Numbers: A Physicalist's Philosophy of Mathematics}. Clarendon Press.
\bibitem[Kukso(2006)]{Kukso2006} Kukso, B. (2006). The reality of absences. \emph{Australasian Journal of Philosophy}, 84(1).
\bibitem[Lewis(2001)]{Lewis2001} Lewis, D. (2001). Truthmaking and difference-making. \emph{No\^us}, 35(4), 602--615.
\bibitem[Lewis(2004)]{Lewis2004} Lewis, D. (2004). Void and object. In J. Collins, N. Hall and L.~A. Paul (Eds.), \emph{Causation and Counterfactuals}. MIT Press.
\bibitem[Martin(1996)]{Martin1996} Martin, C.~B. (1996). How it is: Entities, absences and voids. \emph{Australasian Journal of Philosophy}, 74(1), 57--65.
\bibitem[Molnar(2000)]{Molnar2000} Molnar, G. (2000). Truthmakers for negative truths. \emph{Australasian Journal of Philosophy}, 78(1), 72--86.
\bibitem[Sartre(1943)]{Sartre1943} Sartre, J.-P. (1943). \emph{L'\^Etre et le n\'eant}. Gallimard.
\end{thebibliography}

\end{document}
